% Figure: two-ray ground-reflection received power versus range, showing
% the interference lobing at short range and the R^-4 roll-off far out.
% Relative power in dB for a flat-earth two-ray model, normalised so that
% the free-space level passes through 0 dB at 1 km. The curve is
% P/P_fs = |1 + Gamma exp(j*dphi)|^2 with Gamma = -1 (grazing) and
% dphi = 4*pi*ht*hr/(lambda*R); at large R the sum tends to R^-2 extra
% roll-off, turning the free-space R^-2 into the two-ray R^-4.
\begin{figure}[htbp]
\centering
\begin{tikzpicture}
\begin{axis}[
  width=12cm, height=6.5cm,
  xlabel={range $R$ (km, log scale)},
  ylabel={received power (dB, rel.\ 1\,km free space)},
  xmode=log,
  xmin=0.05, xmax=20, ymin=-40, ymax=25,
  axis lines=left,
  legend pos=south west,
  legend cell align=left,
  domain=0.05:20, samples=400,
]
  % free-space reference: 20*log10(1/R) with R in km, = -20 log10 R
  \addplot[enggray, thick, dashed] { -20*log10(x) };
  \addlegendentry{free space ($R^{-2}$)}
  % two-ray: free-space level times |1 + (-1)*exp(j*dphi)|^2
  % dphi(R) = C / R  with C chosen so first null near 0.1 km for the drawn geometry.
  % |1 - exp(j*dphi)|^2 = 2 - 2 cos(dphi) = 4 sin^2(dphi/2)
  % power_dB = -20 log10 R + 10 log10( 4 sin^2( (C/R)/2 ) )
  % Take C = 3.1416 (rad*km) so dphi = pi at R=1 km -> a lobe structure.
  \addplot[engblue, very thick] {
    -20*log10(x) + 10*log10( 4*(sin( deg( (3.1416/x)/2 ) ))^2 + 0.0001 )
  };
  \addlegendentry{two-ray total}
  % far-field R^-4 asymptote: slope -40 dB/decade through a matched point.
  % at R=10 km the two-ray envelope near -20log10(10) + 10log10(C^2/R^2)-ish;
  % draw a guide line of slope -40 anchored to pass near the far lobes.
  \addplot[engred, thick, dotted] { -40*log10(x) + 6 };
  \addlegendentry{$R^{-4}$ asymptote}
\end{axis}
\end{tikzpicture}
\caption[Two-ray ground-reflection lobing and the $R^{-4}$ roll-off]{Received
power against range for a flat-ground two-ray model, normalised to the
free-space level at one kilometre. Close in, the direct and
ground-reflected rays interfere and the power swings through deep nulls
and lobes as their path difference sweeps through multiples of half a
wavelength. Beyond the last lobe the two rays nearly cancel and the
envelope settles onto an $R^{-4}$ slope, twice as steep as the
free-space $R^{-2}$, the reason a low-mounted terrestrial link loses
range far faster than the Friis equation alone predicts.}
\label{fig:vol04ch10:two-ray-lobing}
\end{figure}
